% =====================================================================
%  ch05b.tex —— 5.2 代表等价的对象与自等价群 + 5.3 半齐次向量丛
%  原书 p.581–p.589 开头（止于 Bibliography 之前；条目由 Lead 在 refs.tex 统一排）
%  本块无二维交换图；编号公式仅 (52)(53)(54)(55)
% =====================================================================

% [原文 p.581]
\section{代表等价的对象与自等价群}

\takeaway{本节求同态 $\gamma_A:\Auteq\Db(A)\to U(A\times\hat A)$ 的核。先证一个独立有用的断言：定义等价的积上对象，在平移意义下其实是层（引理 5.2.1、命题 5.2.2）。由此定出核恰为平移与张量次数零线丛给出的子群 $\Z\oplus(A\times\hat A)_k$，并得到核对象在 $f$ 固定时的唯一性。}

由命题 5.1.12 与命题 5.1.15 可知，存在从正合自等价构成的群 $\Auteq\Db(A)$ 到等距自同构构成的群 $U(A\times\hat A)$ 的同态。本节描述这个同态的核。

由命题 5.1.11 我们知道，所有等价 $\Phi_{(a,\alpha)}[n]$ 都属于这个核。下面证明，核恰好由这些等价构成。为此需要一个本身就有独立意义的断言：对阿贝尔簇，若形如 $\Phi_{\cE}$ 的函子是等价，则积上的对象 $\cE$ 在导出范畴中相差一个平移的意义下其实是一个层。注意，例如对 K3 曲面，这一断言并不成立。

\begin{lemma}{5.2.1}
设 $\cE$ 是 $A\times B$ 上定义等价 $\Phi_{\cE}:\Db(A)\to\Db(B)$ 的对象。考虑投影 $q:(A\times\hat A)\times(B\times\hat B)\to A\times B$，并把直接像 $\Rb q_*\cJ(\cE)$ 记作 $K$，其中 $\cJ(\cE)$ 是 5.1.6 中定义的对象。则 $K$ 同构于对象 $\cE\otimes(\cE^\vee|_{(0,0)})$，这里 $\cE^\vee|_{(0,0)}$ 表示向量空间复形，即对象 $\Rb\cH\mathrm{om}(\cE,\cO_{A\times B})$ 在点 $(0,0)$ 嵌入阿贝尔簇 $A\times B$ 之下的逆像。
\end{lemma}

% [Tutor Note: 文字层把 $\mathbf{R}^{\bullet}\mathcal{H}om$ 抽成 “R·Hom”，页图印作 $\Rb\cH\mathrm{om}$；本译本从页图。]

\begin{proof}
考虑阿贝尔簇
\begin{equation*}
Z=(A\times\hat A)\times(A\times A)\times(B\times B)\times(B\times\hat B)
\end{equation*}
与对象
\begin{equation*}
H=p^*_{1234}\cS_A\otimes p^*_{35}\cE^\vee[n]\otimes p^*_{46}\cE\otimes p^*_{5678}\cS_B^\vee[2n].
\end{equation*}
由关于函子复合的命题 2.1.2 以及图 (51) 可知，$\cJ(\cE)\cong p_{1278*}H$，因此对象 $K$ 等于 $p_{17*}H$。为计算后一对象，先考虑 $Z$ 到
\begin{equation*}
V=A\times(A\times A)\times(B\times B)\times B
\end{equation*}
的投影，并记之为 $v$。为计算 $v_*H$，我们回忆：函子 $\Phi_{\cS_A}$ 是 $\Phi_{\cP_A}$ 与 $\Rf\mu_{A*}$ 的复合，其中
\begin{equation*}
\cP_A=p^*_{14}\cO_\Delta\otimes p^*_{23}\cP\in\Db\bigl((A\times\hat A)\times(A\times A)\bigr).
\end{equation*}
容易看出 $p_{134*}\cP_A\cong\cO_{T_A}[-n]$，其中 $T\subset A\times A\times A$ 是等于 $A$、由点 $(a,0,a)$ 构成的子簇。其次，注意到等式 $\mu_A(a_1,a_2)=(a_1,m(a_1,a_2))$，可以看出 $p_{134*}\cS_A$ 也同构于 $\cO_{T_A}[-n]$。等式 $p_{134*}\cS_B^\vee[2n]=\cO_{T_B}$ 同法验证。

于是，在 $V$ 上有
\begin{equation*}
v_*H\cong p^*_{123}\cO_{T_A}\otimes p^*_{24}\cE^\vee\otimes p^*_{35}\cE\otimes p^*_{456}\cO_{T_B}.
\end{equation*}
考虑嵌入 $j:A\times A\times B\times B\to V$，它由 $(a_1,a_2,b_1,b_2)\mapsto(a_1,0,a_2,0,b_1,b_2)$ 给出。对象 $v_*H$ 同构于 $j_*\cM$，其中
\begin{equation*}
\cM=(\cE^\vee|_{(0,0)})\otimes p^*_{12}\cO_{\Delta_A}\otimes p^*_{23}\cE\otimes p^*_{34}\cO_{\Delta_B}.
\end{equation*}
最后看出 $K\cong p_{14*}\cM\cong(\cE^\vee|_{(0,0)})\otimes\cE$。
\end{proof}

% [原文 p.582]
\begin{proposition}{5.2.2}
设 $A$、$B$ 是阿贝尔簇，$\cE$ 是 $\Db(A\times B)$ 的对象，使得函子 $\Phi_{\cE}:\Db(A)\iso\Db(B)$ 是正合等价。则 $\cE$ 只有一个非平凡上同调层；也就是说，它同构于对象 $\cF[n]$，其中 $\cF$ 是 $A\times B$ 上的层。
\end{proposition}

\yizhu{“在平移意义下是层”指：等价 $\Phi_{\cE}$ 的核对象 $\cE$ 至多相差一次平移 $[n]$ 就是一个层。命题 5.2.2 把这件事归结为几何：$\cJ(\cE)$ 支撑在阿贝尔子簇 $X$ 上且等于 $i_*(L)$（定理 5.1.7），于是 $K=\Rb q'_*L$ 的上同调只落在长度不超过 $d=\dim\Ker q'$ 的区间内；再用 $K\cong\cE^\vee|_{(0,0)}\otimes\cE$ 逐个比较 $\cE$ 各上同调层的指标，就逼出只有一个非零层。原书特别指出，这一断言对 K3 曲面不成立。【来源：原书 p.581–582（断言、证明与“对 K3 曲面不成立”的说明）】}

\begin{proof}
考虑投影 $q:(A\times\hat A)\times(B\times\hat B)\to A\times B$，把它的限制记作 $q'$，即限制到阿贝尔子簇 $X$ 上，$X$ 是层 $\cJ(\cE)$ 的支撑，也是同构 $f_\cE$ 的图。由定理 5.1.7，$\cJ(\cE)$ 等于 $i_*(L)$，其中 $L$ 是 $X$ 上的线丛。把对象 $\Rb q_*\cJ(\cE)=\Rb q'_*L$ 记作 $K$。态射 $q'$ 是阿贝尔簇的同态；令 $d=\dim\Ker(q')$。则 $\dim\im(q')=2n-d$，从而当 $j\notin[0,d]$ 时上同调层 $H^j(K)$ 为零。

另一方面，由引理 5.2.1，$K$ 同构于 $\cE\otimes(\cE^\vee|_{(0,0)})$。

必要时在导出范畴中平移 $\cE$，可以假定 $\cE$ 的最右非零上同调层是 $H^0(\cE)$。设 $H^{-i}(\cE)$（$i\geqslant0$）是 $\cE$ 的最左非零上同调层，$H^k(\cE^\vee)$ 是 $\cE^\vee$ 的最高非零上同调层。必要时用 $T^*_{(a,b)}\cE$ 替换 $\cE$，可以假定点 $(0,0)$ 属于 $H^k(\cE^\vee)$ 的支撑。由于 $\cE$ 的支撑与 $K$ 的支撑重合，可知 $\cE$ 的所有上同调层的支撑都含于 $\im(q')$。特别地，有不等式 $\codim\Supp H^{-i}(\cE)\geqslant d$。因此，对象 $(H^{-i}(\cE))^\vee[-i]$ 的度数小于 $i+d$ 的上同调层为零。

典范态射 $H^{-i}(\cE)[i]\to\cE$ 诱导出非平凡态射
\begin{equation*}
\cE^\vee\longrightarrow(H^{-i}(\cE))^\vee[-i].
\end{equation*}
由于第二个对象的各非平凡上同调层的指标都落在射线 $[i+d,\infty)$ 内，我们看出 $k\geqslant i+d$；这里如前，$H^k(\cE^\vee)$ 是 $\cE^\vee$ 的最高非零上同调层。于是对象
\begin{equation}
K=\cE^\vee|_{(0,0)}\otimes\cE\tag{52}
\end{equation}
具有指标同为 $k\geqslant i+d$ 的非平凡上同调层。另一方面，我们已经知道，当 $j\notin[0,d]$ 时所有上同调层 $H^j(K)$ 为零。这只可能在 $i=0$ 时发生。因此，对象 $\cE$ 只有一个非平凡上同调层，其指标为 $0$，从而它同构于一个层。
\end{proof}

下面考虑 $B\cong A$ 的情形。设 $\cE$ 是 $A\times A$ 上的层，且 $\Phi_{\cE}$ 是自等价。我们要描述所有使 $f_\cE$ 为恒等映射的层 $\cE$；也就是说，其图 $X$ 是 $(A\times\hat A)\times(A\times\hat A)$ 中的对角线。于是对象
\begin{equation*}
K=\cE^\vee|_{(0,0)}\otimes\cE=\Rb q_*\cJ(\cE)
\end{equation*}
具有形式 $\Delta_*(\cM)$，其中 $\cM$ 是 $A$ 上的对象，$\Delta:A\to A\times A$ 是对角嵌入。

% [原文 p.583]
先假定 $(0,0)$ 属于 $\cE$ 的支撑。于是 $\cE^\vee|_{(0,0)}$ 是非平凡的向量空间复形。此时条件 $K=\Delta_*(\cM)$ 蕴涵存在 $A$ 上的层 $E$，使 $\cE\cong\Delta_*(E)$。于是 $\Phi_{\cE}(\,\cdot\,)\cong E\otimes(\,\cdot\,)$。由于 $\Phi_{\cE}$ 是自等价，$E$ 是线丛。容易看出，条件 $f_\cE=\id$ 只能在 $E\in\Pic^0(A)$ 时成立。

若 $(0,0)$ 不属于 $\Supp\cE$，则用层 $\cE':=T_{(a_1,a_2)*}\cE$ 替换 $\cE$，使其支撑包含 $(0,0)$。由命题 5.1.11，$f_{\cE'}=f_\cE$。如上所证，有同构 $\cE'\cong\Delta_*(E')$，其中 $E'\in\Pic^0(A)$。于是 $\cE\cong T_{(a_1-a_2,0)*}\Delta_*(E')$。这就得到下面的推论。

\begin{proposition}{5.2.4}
设 $A$ 是阿贝尔簇。同态 $\gamma_A:\Auteq\Db(A)\to U(A\times\hat A)$ 的核由形如
\begin{equation*}
\Phi_{(a,\alpha)}[i]=T_{a*}(\,\cdot\,)\otimes\cP_\alpha[i]
\end{equation*}
的自等价构成，因而同构于群 $\Z\oplus(A\times\hat A)_k$，其中 $(A\times\hat A)_k$ 是阿贝尔簇 $A\times\hat A$ 的 $k$-值点群。
\end{proposition}

% [Tutor Note: 页图此处印作 “Proposition 5.2.4”，而紧随其后的推论也印作 “Corollary 5.2.4”，编号重复。按 5.2.1、5.2.2 的顺序，本命题应为 5.2.3（本书引言亦引作“命题 5.2.3”）。为不改动原书面貌，此处照录原书编号，请读者注意。]

\begin{corollary}{5.2.4}
设 $A$、$B$ 是两个阿贝尔簇，$\cE_1$、$\cE_2$ 是积 $A\times B$ 上定义凝聚层导出范畴之间等价的对象。此时若 $f_{\cE_1}=f_{\cE_2}$，则
\begin{equation*}
\cE_2\cong T_{a*}\cE_1\otimes\cP_\alpha[i]
\end{equation*}
对某个 $k$-值点 $(a,\alpha)\in A\times\hat A$ 成立。
\end{corollary}

\section{半齐次向量丛}

\takeaway{本节在 $k$ 代数闭、$\operatorname{char}k=0$ 下引入半齐次丛：平移后与自身至多差一个次数零线丛的向量丛。借 Mukai [30] 的滤链、斜率与对应 $\Phi_\mu$，把任意等距同构 $f:A\times\hat A\to B\times\hat B$ 实现成核对象，给出定理 5.1.16 的第二证明，并把 $\Auteq\Db(A)$ 写成一条正合列。}

前几节已经说明，从 $\Db(A)$ 到 $\Db(B)$ 的等价 $\Phi_{\cE}$ 诱导出簇 $A\times\hat A$ 与 $B\times\hat B$ 之间的等距同构。本节假定基域 $k$ 代数闭且 $\operatorname{char}(k)=0$。在这些假设下，我们利用 [30] 的技巧与 [7] 的结果，说明每个等距同构 $f:A\times\hat A\to B\times\hat B$ 都可以按这种方式实现。至于簇 $A\times\hat A$ 与 $B\times\hat B$ 之间存在等距同构蕴涵导出范畴 $\Db(A)$ 与 $\Db(B)$ 等价这一事实，已在 [38] 中证明。因此，我们将对这一事实再给出一个证明。

首先回忆，阿贝尔簇 $D$ 上的每个线丛 $L$ 都给出一个映射 $\varphi_L:D\to\hat D$，它把点 $d$ 映到对应于丛 $T_d^*L\otimes L^{-1}\in\Pic^0(D)$ 的点。这一对应定义了 $\mathrm{NS}(D)$ 到 $\Hom(D,\hat D)$ 的嵌入。此外，众所周知，映射 $\varphi:D\to\hat D$ 属于 $\mathrm{NS}(D)$ 的像，当且仅当 $\hat\varphi=\varphi$。

\yizhu{$\varphi_L$ 称为由线丛 $L$ 给出的同态：$T_d^*L\otimes L^{-1}$ 次数为零，故对应 $\hat D=\Pic^0(D)$ 的一个点；$L$ 丰富时 $\varphi_L$ 是同源，一般只是同态。嵌入 $\mathrm{NS}(D)\hookrightarrow\Hom(D,\hat D)$ 把除子类送到 $\varphi_L$，其像由自伴条件 $\hat\varphi=\varphi$ 刻画；下一段正是用 $\mathrm{NS}(D)\otimes\Q$ 来参数化 $D\times\hat D$ 上的对应。【来源：原书 p.583 正文；Mumford, \emph{Abelian varieties}, Oxford Univ. Press, 1970（本书文献 [32]）】}

阿贝尔簇上的半齐次丛使我们能如下推广上述现象。对 $\mathrm{NS}(D)\otimes\Q$ 的每个元素，都配以 $D\times\hat D$ 上的一个对应；而每个这样的对应都可由某个半齐次丛得到（见下面的命题 5.3.6 与引理 5.1.10）。

先回忆阿贝尔簇上齐次丛与半齐次丛的定义以及若干有关事实。

\begin{definition}{5.3.1}
阿贝尔簇 $D$ 上的向量丛 $\cE$ 称为\emph{齐次}的，如果对每个点 $d\in D$ 都有 $T^*_d(\cE)\cong\cE$。
\end{definition}

% [原文 p.584]
\begin{definition}{5.3.2}
称阿贝尔簇 $D$ 上的向量丛 $\cF$ 是\emph{幂单}的，如果存在滤链
\begin{equation*}
0=\cF_0\subset\cF_1\subset\cdots\subset\cF_n=\cF,
\end{equation*}
使得对一切 $i=1,\dots,n$ 都有 $\cF_i/\cF_{i-1}\cong\cO_D$。
\end{definition}

下面的命题刻画齐次向量丛。

\begin{proposition}{5.3.3}
（[28]、[30]）设 $\cE$ 是阿贝尔簇 $D$ 上的向量丛。则下列条件等价：
\begin{enumerate}[label=(\roman*),leftmargin=2.8em,itemsep=2pt,topsep=2pt]
\item $\cE$ 是齐次的；
\item 存在线丛 $\cP_i\in\Pic^0(D)$ 与幂单丛 $\cF_i$，使得 $\cE\cong\bigoplus_i(\cF_i\otimes\cP_i)$。
\end{enumerate}
\end{proposition}

\begin{definition}{5.3.4}
称阿贝尔簇 $D$ 上的向量丛 $\cE$ 是\emph{半齐次}的，如果对每个点 $d\in D$ 都存在 $D$ 上的线丛 $L$，使得 $T^*_d(\cE)\cong\cE\otimes L$。（我们指出，此时丛 $L$ 属于 $\Pic^0(D)$。）
\end{definition}

\yizhu{半齐次性把齐次性放松为“平移后至多差一个次数零的线丛”：齐次丛要求 $T^*_d\cE\cong\cE$，半齐次丛只要求 $T^*_d\cE\cong\cE\otimes L_d$ 且 $L_d\in\Pic^0(D)$。于是半齐次丛带有两个不变量：斜率 $\mu(\cE)$，以及 $\Pic^0$ 方向上的“平移缺陷”，后者恰好由对应 $\Phi_\mu\subset D\times\hat D$ 记录（命题 5.3.6）。命题 5.3.3 则说明齐次丛正是幂单丛被 $\Pic^0(D)$ 中线丛扭折后的直和。【来源：原书 定义 5.3.4、命题 5.3.3–5.3.6；[30]（Mukai, J. Math. Kyoto Univ. 18 (1978) 239–272）】}

回忆一下：簇上的向量丛称为\emph{单}的，如果它的自同态代数与域 $k$ 重合。

下面的断言在 [30] 中证明。

\begin{proposition}{5.3.5}
（[30]，定理 5.8）设 $\cE$ 是阿贝尔簇 $D$ 上的单向量丛。则下列条件等价：
\begin{enumerate}[label=(\arabic*),leftmargin=2.8em,itemsep=2pt,topsep=2pt]
\item 对任意 $j=0,\dots,n$ 都有 $\dim H^j(D,\cE\mathrm{nd}(\cE))=\binom{n}{j}$；
\item $\cE$ 是半齐次丛；
\item $\cE\mathrm{nd}(\cE)$ 是齐次丛；
\item 存在同源 $\pi:Y\to D$ 与 $Y$ 上的线丛 $L$，使得 $\cE\cong\pi_*(L)$。
\end{enumerate}
\end{proposition}

设 $\cE$ 是阿贝尔簇 $D$ 上的向量丛。把 $\det(\cE)/r(\cE)$ 在 $\mathrm{NS}(D)\otimes_\Z\Q$ 中的等价类记作 $\mu(\cE)$。对 $\mathrm{NS}(D)\otimes_\Z\Q$ 的每个元素 $\mu=[L]/l$，从而对每个丛 $\cE$，都可以配以 $D\times\hat D$ 中的一个对应
\begin{equation*}
\Phi_\mu=\im\bigl[D\xrightarrow{(l,\varphi_L)}D\times\hat D\bigr],
\end{equation*}
其中 $\varphi_L$ 是众所周知的从 $D$ 到 $\hat D$ 的映射，它把 $d$ 映到对应于丛 $T_d^*L\otimes L^{-1}\in\Pic^0(D)$ 的点。若丛是线丛 $L$，就得到映射 $\varphi_L:D\to\hat D$ 的图。把 $\Phi_\mu$ 到 $D$ 与到 $\hat D$ 的投影分别记作 $q_1$ 与 $q_2$。

论文 [30] 给出了所有单半齐次丛的完整描述。

\begin{proposition}{5.3.6}
（[30]，7.10）设 $\mu=[L]/l$，其中 $[L]$ 是丛 $L$ 在 $\mathrm{NS}(D)$ 中的等价类，$l$ 是正整数。则
\begin{enumerate}[label=(\arabic*),leftmargin=2.8em,itemsep=2pt,topsep=2pt]
\item 存在斜率为 $\mu(\cE)=\mu$ 的单半齐次向量丛 $\cE$；
\item 每个斜率为 $\mu(\cE')=\mu$ 的单半齐次向量丛 $\cE'$ 都具有形式 $\cE\otimes M$，其中 $M\in\Pic^0(D)$ 是某个线丛；
\item 有等式 $r(\cE)^2=\deg(q_1)$ 与 $\chi(\cE)^2=\deg(q_2)$。
\end{enumerate}
\end{proposition}

下面的断言使人能用单丛来刻画所有半齐次向量丛。

% [原文 p.585]
\begin{proposition}{5.3.7}
（[30]，6.15、6.16）每个斜率为 $\mu$ 的半齐次向量丛 $\cF$ 都有滤链
\begin{equation*}
0=\cF_0\subset\cF_1\subset\cdots\subset\cF_n=\cF,
\end{equation*}
使得 $\cE_i=\cF_i/\cF_{i-1}$ 都是斜率同为 $\mu$ 的单半齐次向量丛。每个单半齐次丛都是稳定的。
\end{proposition}

% [Tutor Note: 本命题的陈述用 $\cF$ 表示被滤链的丛、用 $\cE_i$ 表示其子商，与定义 5.3.2 用 $\cF_i$ 表示滤链各项的记号不一致；这是原书自身的记号，照录。]

下面两个关于半齐次丛的引理是上述断言的直接推论，在后文中将要用到。

\begin{lemma}{5.3.8}
两个斜率同为 $\mu$ 的单半齐次丛 $\cE_1$ 与 $\cE_2$ 或者同构，或者彼此正交；也就是说，或者 $\cE_1=\cE_2$，或者对一切 $i$ 都有
\begin{equation*}
\Ext^i(\cE_1,\cE_2)=0\quad\text{与}\quad\Ext^i(\cE_2,\cE_1)=0.
\end{equation*}
\end{lemma}

\begin{proof}
由命题 5.3.6，$\cE_2\cong\cE_1\otimes M$，因此 $\cH\mathrm{om}(\cE_1,\cE_2)$ 是齐次丛。由命题 5.3.3，每个齐次丛都可以表示成幂单丛被 $\Pic^0(D)$ 中线丛扭折后的直和。因此，或者 $\cH\mathrm{om}(\cE_1,\cE_2)$ 的所有上同调都消没，从而 $\cE_1$ 与 $\cE_2$ 正交；或者 $\cH\mathrm{om}(\cE_1,\cE_2)$ 有非零截面。后一情形下得到从 $\cE_1$ 到 $\cE_2$ 的非零同态。然而这两个丛都稳定且斜率相同。因此，每个非零同态实际上都是同构。
\end{proof}

\begin{lemma}{5.3.9}
设 $\cE$ 是阿贝尔簇 $D$ 上的单半齐次向量丛。则 $T^*_d(\cE)\cong\cE\otimes\cP_\delta$ 当且仅当 $(d,\delta)\in\Phi_\mu$。
\end{lemma}

\begin{proof}
先说明，对 $\Phi_\mu$ 的任一点 $(d,\delta)$ 都有同构 $T^*_d(\cE)\cong\cE\otimes\cP_\delta$。事实上，令 $l=r(\cE)$、$L=\det(\cE)$。由 $\Phi_\mu$ 的定义，可以把 $(d,\delta)$ 写成 $(lx,\varphi_L(x))$，其中 $x\in D$ 是某个点。由于 $\cE$ 半齐次，存在线丛 $M\in\Pic^0(D)$ 使得
\begin{equation}
T^*_x(\cE)\cong\cE\otimes M.\tag{53}
\end{equation}
两边取行列式，得到等式 $T^*_x(L)\cong L\otimes M^{\otimes l}$。由映射 $\varphi_L$ 的定义，这意味着 $\cP_{\varphi_L(x)}=M^{\otimes l}$。另一方面，把等式 (53) 迭代 $l$ 次，得到
\begin{equation*}
T^*_{lx}(\cE)\cong\cE\otimes M^{\otimes l}=\cE\otimes\cP_{\varphi_L(x)}.
\end{equation*}
于是由 $(d,\delta)=(lx,\varphi_L(x))$ 得到 $T^*_d(\cE)\cong\cE\otimes\cP_\delta$。

现在证逆命题。引入由条件
\begin{equation}
\Sigma^0(\cE):=\{\delta\in\hat D\mid\cE\otimes\cP_\delta\cong\cE\}\tag{54}
\end{equation}
给出的子群 $\Sigma^0(\cE)\subset\hat D$。由于 $\cE$ 半齐次，由命题 5.3.5，$\cE\mathrm{nd}(\cE)$ 是齐次的。于是 $\cE\mathrm{nd}(\cE)$ 可以表示成直和 $\bigoplus_i(\cF_i\otimes\cP_i)$，其中所有 $\cF_i$ 都是幂单的。因此，至多对 $r^2$ 个线丛 $\cP\in\Pic^0(D)$ 有 $H^0(\cE\mathrm{nd}(\cE)\otimes\cP)\neq0$。也就是说，群 $\Sigma^0(\cE)$ 的阶不超过 $r^2$。另一方面，已知
% [原文 p.586]
$q_2(\Ker(q_1))\subset\Sigma^0(\cE)$。于是得到等式 $\operatorname{ord}\Sigma^0(\cE)=r^2$ 与 $q_2(\Ker(q_1))=\Sigma^0(\cE)$。
\end{proof}

\yizhu{$\Sigma^0(\cE)$ 是 $\cE$ 在 $\hat D=\Pic^0(D)$ 方向上的稳定化子：张量次数零线丛后仍同构于自身。引理 5.3.9 把它与对应 $\Phi_\mu$ 的竖直核联系起来：$\Sigma^0(\cE)=q_2(\Ker q_1)$，而由命题 5.3.6 的 (3)，$\deg q_1=r(\cE)^2$、$\deg q_2=\chi(\cE)^2$，两边计数便给出 $\operatorname{ord}\Sigma^0(\cE)=r^2$。“稳定化子的阶等于秩的平方”是后文判定两个纤维丛正交的关键。【来源：原书 式 (54)、引理 5.3.9、命题 5.3.6；[30] §7】}

% [Tutor Note: “另一方面，已知”一句原文只写 “it is known”，未给证明或出处，属跳步；按命题 5.3.6 的 (3) 与 $\Ker q_1$ 的构造可以补齐（译者核对）。]

现在假定对某个点 $(d,\delta)\in D\times\hat D$ 有 $T^*_d(\cE)\cong\cE\otimes\cP_\delta$。考虑满足 $(d,\delta')\in\Phi_\mu$ 的点 $\delta'\in\hat D$。如上所证，此时有同构 $T^*_d(\cE)\cong\cE\otimes\cP_{\delta'}$。于是 $\cE\otimes\cP_{(\delta-\delta')}\cong\cE$，从而 $\delta-\delta'\in\Sigma^0(\cE)$。然而，由于 $\Sigma^0(\cE)=q_2(\Ker(q_1))$，可知点 $(0,\delta-\delta')$ 属于 $\Phi_\mu$。因此，点 $(d,\delta)$ 也属于 $\Phi_\mu$。

设 $f$ 是等距同构。下面给出一项构造，说明如何由 $f$ 构造积上的对象 $\cE$，使得 $\cE$ 定义导出范畴的等价，并且 $f_\cE$ 与 $f$ 重合。

\textbf{构造~5.3.10}\ 固定等距同构 $f:A\times\hat A\to B\times\hat B$，把它的图记作 $\Gamma$。如上，把同构 $f$ 写成矩阵形式
\begin{equation*}
f=\begin{pmatrix}x&y\\ z&w\end{pmatrix}.
\end{equation*}
假定 $y:\hat A\to B$ 是同源。此时对映射 $f$ 配以元素 $g\in\Hom(A\times B,\hat A\times\hat B)\otimes_\Z\Q$，其形式为
\begin{equation*}
g=\begin{pmatrix}y^{-1}x&-y^{-1}\\ -\hat y^{-1}&wy^{-1}\end{pmatrix}.
\end{equation*}
元素 $g$ 定义了 $(A\times B)\times(\hat A\times\hat B)$ 上的某个对应。容易看出，$f$ 是等距的这一条件蕴涵等式 $\hat g=g$。这意味着 $g$ 实际上属于 $\mathrm{NS}(A\times B)\otimes_\Z\Q$ 在到 $\Hom(A\times B,\hat A\times\hat B)\otimes_\Z\Q$ 的典范嵌入下的像（例如见 [32]）。于是存在 $\mu=[L]/l\in\mathrm{NS}(A\times B)$，使得 $\Phi_\mu$ 与对应 $g$ 的图重合。命题 5.3.6 告诉我们，对任一 $\mu$ 都可以构造 $A\times B$ 上斜率为 $\mu(\cE)=\mu$ 的单半齐次丛 $\cE$。

% [Tutor Note: 原书此处用独立的 “Construction 5.3.10” 环境；本书陈述环境表中没有 construction，故按体例“原文不用标准环境者作普通段落”处理，以粗体引出。另：$g$ 的矩阵项 $-\hat y^{-1}$ 指对偶同源 $\hat y:\hat B\to\hat A$ 之逆，原书未加说明，按上下文补出。]

下面将说明从 $\Db(A)$ 到 $\Db(B)$ 的函子 $\Phi_{\cE}$ 是等价，且 $f_\cE=f$。不过先把两个图 $\Gamma$ 与 $\Phi_\mu$ 作比较。若点 $(a,\alpha,b,\beta)$ 属于 $\Gamma$，则
\begin{equation*}
b=x(a)+y(\alpha),\qquad \beta=z(a)+w(\alpha),
\end{equation*}
从而
\begin{equation*}
\alpha=-y^{-1}x(a)+y^{-1}(b),\qquad \beta=(z-wy^{-1}x)(a)+wy^{-1}(b).
\end{equation*}
由于 $f$ 是等距的，有等式 $(z-wy^{-1}x)=-\hat y^{-1}$。因此，点 $(a,\alpha,b,\beta)$ 属于图 $\Gamma$ 当且仅当 $(a,-\alpha,b,\beta)$ 属于 $\Phi_\mu$。所以
\begin{equation*}
\Phi_\mu=(1_A,-1_{\hat A},1_B,1_{\hat B})\Gamma.
\end{equation*}
特别地，由于 $f$ 是同构，可知 $\Phi_\mu$ 到 $A\times\hat A$ 与 $B\times\hat B$ 的投影都是同构。

% [原文 p.587]
\begin{proposition}{5.3.11}
设 $\cE$ 是 $A\times B$ 上由如上所述的等距同构 $f$ 构造出的半齐次丛。则函子 $\Phi_{\cE}:\Db(A)\to\Db(B)$ 是等价。
\end{proposition}

\begin{proof}
把 $\cE$ 在纤维 $\{a\}\times B$ 上的限制记作 $\cE_a$。由定理 2.1.5，为证明函子 $\Phi_{\cE}$ 全忠实，只需说明所有丛 $\cE_a$ 都是单的，且对不同点的丛彼此正交。

首先注意，由命题 5.3.6，$\cE$ 的秩等于映射 $\Phi_\mu\to A\times B$ 的次数的平方根，即 $\sqrt{\deg(\beta)}$。

由于 $\cE$ 半齐次，立即可知所有丛 $\cE_a$ 也都半齐次。此外，限制的斜率 $\mu(\cE_a)$ 等于 $\delta\beta^{-1}\in\mathrm{NS}(B)\otimes\Q\subset\Hom(B,\hat B)\otimes\Q$。为简便起见，把 $\delta\beta^{-1}$ 记作 $\nu$，并将它看作 $\mathrm{NS}(B)\otimes\Q$ 的元素。由命题 5.3.6，存在 $B$ 上斜率为给定值 $\mu(\cF)=\nu$ 的单半齐次丛 $\cF$。显然，此时映射 $\Phi_\nu$ 就是 $\im\bigl[\hat A\xrightarrow{(\beta,\delta)}B\times\hat B\bigr]$。由于 $f$ 是同构，可知 $\hat A\xrightarrow{(\beta,\delta)}B\times\hat B$ 是嵌入。于是再次应用命题 5.3.6，得到等式 $r(\cF)=\sqrt{\deg(\beta)}=r(\cE_a)$。这样，两个丛 $\cF$ 与 $\cE_a$ 半齐次，且斜率相同、秩相同。此外，丛 $\cF$ 是单的。由命题 5.3.7 与命题 5.3.6 的 (2) 可知，$\cE_a$ 也是单丛。

其次，由引理 5.3.8，对两个点 $a_1,a_2\in A$，丛 $\cE_{a_1}$ 与 $\cE_{a_2}$ 或者正交，或者同构。假设它们同构。由于 $\cE$ 半齐次，可知存在点 $(\alpha,\beta)\in\hat A\times\hat B$ 使
\begin{equation}
T^*_{(a_2-a_1,0)}\cE\cong\cE\otimes\cP_{(\alpha,\beta)}.\tag{55}
\end{equation}
特别地，得到
\begin{equation*}
\cE_{a_2}\otimes\cP_\beta\cong\cE_{a_1}\cong\cE_{a_2}.
\end{equation*}
因此 $\cP_\beta\in\Sigma^0(\cE_a)$（见 (54)）。

由引理 5.3.9 与命题 5.3.6，群 $\Sigma^0(\cE)$ 与 $\Sigma^0(\cE_a)$ 的阶都等于 $r^2$。我们断言自然映射 $\sigma:\Sigma^0(\cE)\to\Sigma^0(\cE_a)$ 是同构。事实上，否则将存在点 $\alpha'\in\hat A$ 使 $\cE\otimes\cP_{\alpha'}\cong\cE$。于是由引理 5.3.9，$(0,\alpha',0,0)\in\Phi_\mu$。这与投影 $\Phi_\mu\to B\times\hat B$ 是同构这一事实矛盾。

现在，若 $\sigma$ 是同构，则存在点 $\alpha'\in\hat A$ 使 $\cE\otimes\cP_{(\alpha',\beta)}\cong\cE$。由 (55) 可知
\begin{equation*}
T^*_{(a_2-a_1,0)}\cE\cong\cE\otimes\cP_{(\alpha-\alpha',0)}.
\end{equation*}
由引理 5.3.9，这意味着点 $(a_2-a_1,\alpha-\alpha',0,0)$ 属于 $\Phi_\mu$。由于投影 $\Phi_\mu\to B\times\hat B$ 是同构，再次得到等式 $a_2-a_1=0$。于是，对两个不同的点 $a_1$ 与 $a_2$，对应的丛 $\cE_{a_1}$ 与 $\cE_{a_2}$ 正交。因此，函子 $\Phi_{\cE}:\Db(A)\to\Db(B)$ 全忠实。同理，伴随函子 $\Psi_{\cE^\vee}$ 也全忠实。因此 $\Phi_{\cE}$ 是等价。
\end{proof}

% [原文 p.588]
\begin{proposition}{5.3.12}
设 $\cE$ 是由如上所述的等距同构 $f:A\times\hat A\to B\times\hat B$ 构造出的半齐次丛。则 $f_\cE=f$。
\end{proposition}

\begin{proof}
把态射 $f_\cE$ 的图记作 $X$。由推论 5.1.10，点 $(a,\alpha,b,\beta)$ 属于 $X$ 当且仅当
\begin{equation*}
T_{b*}\cE\otimes\cP_\beta\cong T^*_a\cE\otimes\cP_\alpha,
\end{equation*}
这等价于等式
\begin{equation*}
T^*_{(a,b)}\cE\cong\cE\otimes\cP_{(-\alpha,\beta)}.
\end{equation*}
于是由引理 5.3.9 看出 $X=(1_A,-1_{\hat A},1_B,1_{\hat B})\Phi_\mu$，其中 $\mu=\mu(\cE)$ 是 $\cE$ 的斜率。另一方面，由构造 5.3.10，$f$ 的图 $\Gamma$ 也等于 $(1_A,-1_{\hat A},1_B,1_{\hat B})\Phi_\mu$。因此，同构 $f_\cE$ 与 $f$ 重合。
\end{proof}

在由同构 $f$ 构造丛 $\cE$ 时，我们假定了映射 $y:\hat A\to B$ 是同源。若情形并非如此，则把 $f$ 表示成两个映射 $f_1\in U(A\times\hat A,B\times\hat B)$ 与 $f_2\in U(A\times\hat A)$ 的复合，使 $y_1$ 与 $y_2$ 都是同源。容易看出这总是可以做到的。于是对每个映射 $f_i$ 找出相应的对象 $\cE_i$，考虑函子 $\Phi_{\cE_i}$ 的复合，并取代表它的对象。

本节与前几节证明的断言可以汇集为下面的定理。

\begin{theorem}{5.3.13}
设 $A$、$B$ 是特征 $0$ 的代数闭域上的两个阿贝尔簇。则凝聚层的有界导出范畴 $\Db(A)$ 与 $\Db(B)$ 作为三角范畴等价，当且仅当存在等距同构 $f:A\times\hat A\iso B\times\hat B$。
\end{theorem}

\begin{theorem}{5.3.14}
设 $A$ 是特征 $0$ 的代数闭域上的阿贝尔簇。则导出范畴的正合自等价群 $\Auteq\Db(A)$ 适合如下群的短正合列：
\begin{equation*}
0\longrightarrow\Z\oplus(A\times\hat A)_k\longrightarrow\Auteq\Db(A)\longrightarrow U(A\times\hat A)\longrightarrow 1.
\end{equation*}
\end{theorem}

\yizhu{这条正合列把 $\Auteq\Db(A)$ 分成两层：离散部分 $\Z$ 来自平移 $[i]$，几何部分 $(A\times\hat A)_k$ 来自 $T_{a*}(\,\cdot\,)\otimes\cP_\alpha$；商 $U(A\times\hat A)$ 是 $A\times\hat A$ 的等距自同构群。序列不典范分裂，其障碍即末段所说的 2-上闭链（2-cocycle），显式公式见 [37]。于是“计算自等价群”被拆成两件可分别处理的事：算 $\Z\oplus(A\times\hat A)_k$，以及算中心扩张的类。【来源：原书 p.588 与 [37]（Orlov, Izv. Math. 66 (2002) 569–594）】}

于是，群 $\Auteq\Db(A)$ 有正规子群 $(A\times\hat A)_k$，它由形如 $T_{a*}(\,\cdot\,)\otimes\cP_\alpha$ 的函子构成，其中 $(a,\alpha)\in A\times\hat A$。关于该子群的商是 $U(A\times\hat A)$ 被 $\Z$ 的中心扩张。这个中心扩张由某个 2-上闭链描述，其公式可在 [37] 中找到。

\begin{example}{5.3.15}
考虑自同态环 $\End(A)$ 同构于 $\Z$ 的阿贝尔簇 $A$。则 Néron--Severi 群 $\mathrm{NS}(A)$ 同构于 $\Z$。把 $\mathrm{NS}(A)$ 与 $\mathrm{NS}(\hat A)$ 的生成元分别记作 $L$ 与 $M$。复合 $\varphi_M\circ\varphi_L$ 等于 $N\cdot\id_A$，其中 $N>0$ 是某个整数。此时群 $U(A\times\hat A)$ 与同余子群 $\Gamma_0(N)\subset\mathrm{SL}(2,\Z)$ 重合。其次，设 $B$ 是另一个阿贝尔簇，使 $B\times\hat B$ 同构于 $A\times\hat A$。容易看出，每个这样的同构都是等距的。阿贝尔簇 $B$ 可以表示成某个态射
\begin{equation*}
A\xrightarrow{(k\cdot\id,\,m\varphi_L)}A\times\hat A
\end{equation*}
的像。可以假定 $\gcd(k,m)=1$。

% [原文 p.589]
把这样定义的从 $A$ 到 $B$ 的态射记作 $\psi$。$\psi$ 的核是 $\Ker(m\varphi_L)\cap A_k$。由于 $\gcd(k,m)=1$，实际上有 $\Ker(\psi)=\Ker(\varphi_L)\cap A_k$。另一方面，有包含关系 $\Ker(\varphi)\subset A_N$。于是不失一般性，可以假定 $k$ 是 $N$ 的因子。$N$ 的每个因子 $k$ 都诱导出形如 $B:=A/(\Ker(\varphi_L)\cap A_k)$ 的阿贝尔簇。显然，$N$ 的两个不同因子给出不同构的阿贝尔簇。此外，容易看出，$B$ 在 $A\times\hat A$ 中的嵌入可裂，当且仅当 $\gcd(k,N/k)=1$。因此，使 $\Db(B)\simeq\Db(A)$ 的阿贝尔簇 $B$ 的个数等于 $2^s$，其中 $s$ 是 $N$ 的素因子个数。
\end{example}

% [Tutor Note: 末段中的 $\Ker(\varphi)$ 是 $\Ker(\varphi_L)$ 的简写（前几句均写 $\varphi_L$）；原书如此，照录。]

\chapterreview{%
第 5 章是全书最后一个实质性章节，也是第 3 章主定理最完整的一次应用。第 3 章说，光滑射影簇之间的等价必由积上的核对象给出；第 4 章据此把 K3 曲面导出范畴的等价翻译成 Mukai 格上的霍奇等距；第 5 章对阿贝尔簇给出同样干净的答案，但答案的形状不同：K3 的答案落在格论上，阿贝尔簇的答案落在 $A\times\hat A$ 的自同构群上。之所以能算得这么清楚，是因为阿贝尔簇有群结构：平移提供大量自同构，对偶阿贝尔簇又把 $\Pic^0$ 变成几何对象，于是“等价”被译成“等距同构”这一较易检验的条件。

全章分三步。5.1 建立基本字典：正合自等价诱导 $A\times\hat A$ 的等距同构，得到同态 $\gamma_A:\Auteq\Db(A)\to U(A\times\hat A)$，并证明等价由核对象唯一决定。5.2 做减法：先证明一个独立有用的断言——阿贝尔簇上定义等价的积上对象，在平移意义下其实是层（引理 5.2.1、命题 5.2.2），据此定出 $\gamma_A$ 的核恰为平移与张量次数零线丛构成的子群（命题 5.2.4），并得到 $f_\cE$ 相同时核对象的唯一性（推论 5.2.4）。5.3 做加法：反过来，从任意等距同构 $f$ 出发，用 Mukai [30] 的半齐次丛技术造出核对象（构造 5.3.10、命题 5.3.11、5.3.12），既补上定理 5.1.16 的第二证明，又把自等价群写成一条正合列（定理 5.3.14）。两步之间靠的是同一个对象：$\mathrm{NS}(D)\otimes\Q$ 中的斜率 $\mu$ 与它给出的对应 $\Phi_\mu\subset D\times\hat D$。

结论可以归为三条。其一，$\Db(A)\simeq\Db(B)$ 当且仅当 $A\times\hat A$ 与 $B\times\hat B$ 等距同构（定理 5.1.16，定理 5.3.13）。其二，$\Auteq\Db(A)$ 的核是 $\Z\oplus(A\times\hat A)_k$：$\Z$ 由平移 $[i]$ 生成，$(A\times\hat A)_k$ 由 $T_{a*}(\,\cdot\,)\otimes\cP_\alpha$ 给出；商是等距自同构群，故 $\Auteq\Db(A)$ 是 $U(A\times\hat A)$ 被 $\Z$ 的中心扩张，其类由一个 2-上闭链描述（定理 5.3.14）。其三，核对象在 $f$ 固定时唯一到平移与 $\Pic^0$ 扭折（推论 5.2.4）；例 5.3.15 更把它算到底：当 $\End(A)=\Z$ 时，与 $A$ 导出等价的 $B$ 恰有 $2^s$ 个，$s$ 为 $N$ 的素因子个数。

读这一章有几处需要提防。第一，原书 p.583 把核定理印作“命题 5.2.4”，紧随其后的推论也印作“推论 5.2.4”，编号撞车；按 5.2.1、5.2.2 的顺序它应为命题 5.2.3，本书引言亦引作 5.2.3。译本照录原号并在该处加注。第二，“相差平移是层”这条断言只对阿贝尔簇成立，原书明说对 K3 曲面不成立——这恰是阿贝尔簇比 K3 好算的地方，不可当作一般现象。第三，构造 5.3.10 预先假定 $y:\hat A\to B$ 是同源，一般情形靠把 $f$ 分解成两个满足条件的等距同构之复合；原书只说“容易看出这总是可以做到的”，未给细节。第四，p.585 末“it is known that $q_2(\Ker q_1)\subset\Sigma^0(\cE)$”一句没有证明，须回到命题 5.3.6 的次数计算自行补齐。第五，5.3 节全程要求 $k$ 代数闭且 $\operatorname{char}k=0$。

从方法上看，这一章最值得学的是“两条腿走路”。5.1–5.2 从函子出发，靠同调代数把核压成 $\Z\oplus(A\times\hat A)_k$：先把核对象在平移意义下做成层，再逐层比较上同调指标。5.3 反过来从几何出发，靠半齐次丛把任意等距同构实现成核对象：把 $f$ 翻成 $\mathrm{NS}(A\times B)\otimes\Q$ 中的元素 $g$，再由 $g$ 定出斜率、按命题 5.3.6 造单丛。两条路线在定理 5.3.13 处会合，互为验证。这也说明了半齐次丛为何在此处是“对的工具”：它把 $\mathrm{NS}(D)\otimes\Q$ 的每个元素变成一个 $D\times\hat D$ 上的对应，而等距同构正是这样的对应，于是“造核对象”退化为“按给定斜率造单丛”这一有现成答案的问题。留下的问题也很清楚：正特征、非代数闭域，以及中心扩张那个 2-上闭链的显式公式（原书指向 [37]），都在这套框架的覆盖之外。}
